% Part: second-order-logic
% Chapter: syntax-and-semantics
% Section: expressive-power

\documentclass[../../../include/open-logic-section]{subfiles}

\begin{document}

\olfileid{sol}{syn}{exp}
\olsection{توان بیانی}

\begin{explain}
کمیت‌گذاری بر متغیرهای مرتبهٔ دوم سبب افزایش چشمگیر توان بیانی زبان در
مقایسه با منطق مرتبهٔ اول می‌شود. کمیت‌گذاری وجودی مرتبهٔ دوم به ما
اجازه می‌دهد بگوییم توابع یا روابطی با ویژگی‌های معین وجود دارند. در
منطق مرتبهٔ اول، تنها راه انجام این کار آن است که نمادی غیرمنطقی (یعنی
!!a{function} یا !!{predicate}) برای این مقصود تعیین کنیم. کمیت‌گذاری
کلی مرتبهٔ دوم به ما امکان می‌دهد بگوییم همهٔ زیرمجموعه‌های !!{domain}،
همهٔ روابط بر آن، یا همهٔ توابع از !!{domain} به خودش ویژگی‌ای دارند.
در منطق مرتبهٔ اول فقط می‌توانیم بگوییم زیرمجموعه‌ها، روابط یا توابعی
که به یکی از نمادهای غیرمنطقی زبان نسبت داده شده‌اند ویژگی‌ای دارند.
همچنین هنگامی که می‌گوییم زیرمجموعه‌ها، روابط یا توابعی با یک ویژگی
وجود دارند، یا همهٔ آن‌ها آن ویژگی را دارند، در تعیین خود این ویژگی
نیز می‌توانیم از کمیت‌گذاری مرتبهٔ دوم استفاده کنیم. بدین‌سان می‌توانیم
روابطی تعریف کنیم که در منطق مرتبهٔ اول تعریف‌پذیر نیستند و ویژگی‌هایی
از دامنه را بیان کنیم که در منطق مرتبهٔ اول بیان‌پذیر نیستند.
\end{explain}

\begin{defn}
اگر~$\Struct{M}$ !!a{structure} برای زبان~$\Lang{L}$ باشد، رابطهٔ
$R \subseteq \Domain{M}^2$ در~$\Lang{L}$ \emph{تعریف‌پذیر} است هرگاه
!!{formula}ی چون~$!A_R(x, y)$ وجود داشته باشد که فقط متغیرهای~$x$
و~$y$ در آن آزاد باشند، به‌گونه‌ای که~$R(a, b)$ برقرار باشد (یعنی
$\tuple{a,b} \in R$) اگر و تنها اگر~$\Sat{M}{!A_R(x, y)}[s]$ برای
تخصیصی برقرار باشد که~$s(x) = a$ و~$s(y) = b$.
\end{defn}

\begin{ex}
در منطق مرتبهٔ اول می‌توانیم رابطهٔ همانی~$\Id{\Domain{M}}$ (یعنی
$\Setabs{\tuple{a,a}}{a \in \Domain{M}}$) را با فرمول~$\eq[x][y]$
تعریف کنیم. در منطق مرتبهٔ دوم می‌توانیم این رابطه را
\emph{بدون~$\eq$} تعریف کنیم. زیرا اگر~$a$ و~$b$ همان !!{element}
از~$\Domain{M}$ باشند، آنگاه !!{element}s زیرمجموعه‌های یکسانی از
$\Domain{M}$ هستند (چون مجموعه‌ها به‌وسیلهٔ !!{element}s خود تعیین
می‌شوند). برعکس، اگر~$a$ و~$b$ متفاوت باشند، !!{element}s
زیرمجموعه‌های یکسانی نیستند: برای
نمونه، $a \in \{a\}$ اما~$b \notin \{a\}$، هرگاه~$a \neq b$.
پس «!!{element}s زیرمجموعه‌های یکسانی از~$\Domain{M}$ بودن» رابطه‌ای
است که دربارهٔ~$a$ و~$b$ برقرار است اگر و تنها اگر~$a = b$. این رابطه
را می‌توان در منطق مرتبهٔ دوم بیان کرد، زیرا می‌توانیم بر همهٔ
زیرمجموعه‌های~$\Domain{M}$ کمیت‌گذاری کنیم. بنابراین !!{formula} زیر
$\Id{\Domain{M}}$ را تعریف می‌کند:
\[
\lforall[X][(X(x) \liff X(y))]
\]
\end{ex}

\begin{prob}
نشان دهید~$\lforall[X][(X(x) \lif X(y))]$ (توجه کنید:
$\lif$، نه~$\liff$!) رابطهٔ~$\Id{\Domain{M}}$ را تعریف می‌کند.
\end{prob}

\begin{ex}
اگر~$R$ یک !!{predicate} دوموضعی باشد، $\Assign{R}{M}$ رابطه‌ای
دوموضعی بر~$\Domain{M}$ است. شاید قدری گیج‌کننده باشد، اما~$R$ را هم
برای !!{predicate}~$R$ و هم برای خود رابطهٔ~$\Assign{R}{M}$ به کار
خواهیم برد. \emph{بستار تعدی}~$R^*$ از~$R$ رابطه‌ای است که میان~$a$ و
$b$ برقرار است اگر و تنها اگر برای~$c_1$، \dots، $c_k$ داشته باشیم
$R(a,c_1)$، $R(c_1, c_2)$، \dots، $R(c_k,b)$. این تعریف حالت~$k = 0$
را نیز شامل می‌شود؛ یعنی اگر~$R(a,b)$ برقرار باشد، $R^*(a,b)$ نیز
برقرار است. این بدان معناست که~$R \subseteq R^*$. در واقع، $R^*$
کوچک‌ترین رابطه‌ای است که~$R$ را در بر می‌گیرد و تعدی است. در منطق
مرتبهٔ دوم می‌توانیم بگوییم~$X$ رابطه‌ای تعدی است که~$R$ را در بر
می‌گیرد:
\begin{multline*}
  !B_R(X) \ident \lforall[x][\lforall[y][(R(x,y) \lif X(x, y))]] \land {}\\
\lforall[x][\lforall[y][\lforall[z][((X(x,y) \land X(y,z)) \lif X(x,
      z))]]].
\end{multline*}
هموند نخست می‌گوید~$R \subseteq X$ و هموند دوم می‌گوید~$X$ تعدی است.

گفتن اینکه~$X$ کوچک‌ترین چنین رابطه‌ای است، یعنی خود آن در هر رابطه‌ای
که~$R$ را در بر می‌گیرد و تعدی است گنجانده می‌شود. پس می‌توانیم بستار
تعدی~$R$ را با !!{formula} زیر تعریف کنیم:
\[
R^*(X) \ident !B_R(X) \land \lforall[Y][(!B_R(Y) \lif
  \lforall[x][\lforall[y][(X(x, y) \lif Y(x,y))]])].
\]
داریم~$\Sat{M}{R^*(X)}[s]$ اگر و تنها اگر~$s(X) = R^*$. بستار تعدی
$R$ را نمی‌توان در منطق مرتبهٔ اول بیان کرد.
\end{ex}

\end{document}
